\section{Background and Related Work}
\label{sec:related}

\subsection{SOFL-Oriented Fault Prevention}
Li and Liu studied requirements-related fault prevention in formal-specification-driven development, including dependency-aware ordering and transformation guidance~\cite{li2022qrs,li2023iet}.
These works motivate prevention-first design, but they do not directly address integration with current LLM-assisted iterative workflows.

\subsection{LLM + Formal Verification}
Verifier-guided synthesis has shown encouraging results in Dafny and related domains, commonly using iterative correction and conformance-centric metrics~\cite{lynn2024tosem,dafnybench2024}.
This literature supports the use of explicit feedback loops, but transferability to Agile-SOFL scenario logic remains an open empirical question.

\subsection{Specification Mutation and Fault-Based Testing}
Specification mutation is a long-established method for evaluating fault sensitivity of testing and analysis processes~\cite{okun2007nist,circus2017mutation}.
Following this tradition, we adopt two mutation layers (specification confusion and implementation screening) to assess prevention behavior beyond simple pass/fail outcomes.

\subsection{Counterexample-Guided Inductive Synthesis}
Counterexample-guided inductive synthesis (CEGIS), introduced by Solar-Lezama et al.~\cite{solar2006sketch}, provides the foundational template for iterative correction loops in program synthesis.
In the CEGIS schema, a verifier supplies counterexamples to a learner, which refines its candidate until no counterexample can be found.
HSP-Agile instantiates a practical approximation of this loop: the formal checker plays the role of the verifier, emitting concrete failing cases, while the LLM plays the role of the learner responding to structured feedback.
Because our setting targets Agile-SOFL scenario logic rather than full program proofs, convergence is not guaranteed, but the repair loop consistently reduces conformance violations across the attempt budget.

\subsection{Static Analysis and Pattern-Based Defect Detection}
Static analysis tools such as FindBugs~\cite{hovemeyer2004findbugs} and its successor SpotBugs operate by matching program structures against a catalog of known bug patterns, assigning severity ratings without requiring test execution.
This design---lightweight, oracle-free, and fast---motivates the pattern guard component of HSP-Agile.
Where static analysis tools target implementation-level code smells, our pattern guard encodes requirement-level anti-patterns specific to SOFL scenario logic (e.g., branch-order confusion, missing guard refinement), extending the defect-detection paradigm to the specification-to-code translation phase.

\subsection{Formal Methods in Agile Development}
The tension between formal rigor and agile delivery pace is well-studied in software engineering~\cite{agile2001manifesto}.
Lightweight formal methods---model checking over finite state abstractions, property-based testing, and runtime verification---have been proposed as practical compromises that preserve agile iteration speed while providing bounded correctness guarantees~\cite{clarke1999modelchecking}.
HSP-Agile contributes to this body of work by demonstrating that a formally grounded acceptance predicate, instantiated through scenario-level checking and repair, can be integrated into a per-sprint delivery cadence without prohibitive latency overhead.

\subsection{Gap and Positioning}
Existing approaches often emphasize either correctness guarantees or interactive efficiency.
Our study positions HSP-Agile as a hybrid approach and evaluates it against one-shot and test-feedback baselines under a benchmark designed to expose precedence-sensitive failures.
The goal is comparative empirical evidence rather than universal claims of superiority.
